\section{Materials and Methods}
\label{sec:methods}

\subsection{Study Selection and OSDR Metadata Extraction}
Candidate RNA-seq datasets were identified through systematic query of the NASA Open Science Data Repository (OSDR, \url{https://osdr.nasa.gov}) utilizing criteria established by \citet{barker2023meta}. Eligible accessions satisfied three requirements: (1) experimental cultivation of \textit{Arabidopsis thaliana} seedlings aboard the International Space Station (ISS), (2) inclusion of synchronous or asynchronous ground hardware controls, and (3) availability of raw paired-end or single-end FASTQ sequencing files. 

Six primary accessions were selected for comparative meta-transcriptomic dissection (Table~\ref{tab:datasets}): OSD-37 (GLDS-37, BRIC-19; 8-day-old dark-grown seedlings across Col-0, Cvi-0, Ler-0, WS-2), OSD-38 (GLDS-38, BRIC-20; 3--4-day-old dark-grown Col-0 seedlings with liquid nitrogen preservation), OSD-120 (GLDS-120, APEX-03-2; 11--13-day-old micro-dissected root tips in light/dark cycles), OSD-217 (GLDS-217, BRIC-PDF; light-grown WS seedlings with separated leaf and root tissues), OSD-321 (GLDS-321, European Modular Cultivation System; microgravity vs.\ 1\,$g$ on-orbit centrifuge), and OSD-522 (GLDS-522; light-grown Col-0 seedlings). Study metadata and sequencing run manifests were extracted programmatically using the OSDR RESTful API and assembled into computational tables via automated Python/Google Colab scripts.

\subsection{Quality Trimming and Host-Genome Depletion}
Raw sequencing reads were retrieved in gzipped FASTQ format. Quality verification was performed using FastQC (v0.11.9) \citep{ewels2016multiqc}, and adapter contamination and low-quality bases were filtered using Trim Galore! (v0.6.7) with a Phred score cutoff of $Q \ge 25$ and minimum read length of 35\,bp. 

To deplete host-derived sequences, cleaned reads were aligned against the \textit{Arabidopsis thaliana} reference genome (TAIR10 / Araport11) using the splice-aware aligner STAR (v2.7.9a) and Bowtie2 (v2.4.4). Reads failing to align to the plant nuclear, mitochondrial, or chloroplastic genomes (designated as unmapped FASTQ reads) were extracted using SAMtools (v1.14) with flags \texttt{-f 4} (single-end) or \texttt{-f 12 -F 256} (paired-end), and retained as the non-host metatranscriptome library.

\subsection{Taxonomic Profiling and Abundance Estimation}
Taxonomic classification of the microbial read pool was executed using MetaPhlAn4 (v4.0.6) \citep{segata2012metagenomic}, which relies on a curated database of $\sim$5.1 million clade-specific marker genes covering $>26,000$ microbial species (Bacteria, Archaea, Eukaryota, and Viruses). In parallel, Bowtie2 alignment files were generated to assess mapping densities across marker loci. Relative abundance estimates were calculated at each hierarchical rank (domain, phylum, class, order, family, genus, and species). Interactive hierarchical taxonomic visualizations were generated using Krona (v2.8.1) \citep{ondov2011interactive} to inspect nested clades across experimental conditions.

\subsection{Intersectional and Community Diversity Analysis}
To determine core versus flight-specific or ground-specific taxa, intersectional set analyses were conducted using UpSetR (v1.4.0) \citep{conway2017upsetr} and Intervene (v0.6.5). Pairwise Jaccard similarities and Spearman correlation coefficients were calculated between samples to construct correlation matrices and hierarchical dendrograms. Hierarchical All-against-All Association (HAllA) testing was implemented to identify co-varying microbial clusters across mission parameters and ecotypes.

\subsection{Functional Annotation and Ecological Provenance}
Identified microbial species were annotated against The Microbe Directory (v2.0) and BacDive (The Bacterial Diversity Metadatabase) to assign ecological provenance and functional traits. Microorganisms were systematically classified into: (1) \textit{Putative Seed Endophytes / Plant Associates} (rhizosphere, phyllosphere, or seed-borne symbionts), (2) \textit{Human / Clinical Commensals} (known inhabitants of human skin, oral cavity, or gut), and (3) \textit{Environmental / Spacecraft Hardware Contaminants} (oligotrophic water-system or cleanroom isolates). Physiological attributes—including Gram stain morphology, biofilm-forming capacity, spore formation, optimal growth temperature, and pathogenicity—were compiled to evaluate their functional potential in closed spaceflight habitats.
